% §4 Modelling assumptions.  A1 (GTR) and A2 (clock), the constant block model they induce at the
% root split, and the two scalars eta and rho in which the sampling rate is stated.  C1 is on its
% own counter: it is an inequality on the tree, not a hypothesis on the data.  The proof of
% prop:cfbm is deferred to app:comp1; the Fiedler-vector shape belongs to lem:A-decomp, not here.
% No subsections: assume a clock -> get the block model -> read eta and rho off it is one argument.
\section{Modelling assumptions}\label{sec:assumptions}

In this section we state the two hypotheses the analysis carries, the block model they put the
similarity matrix in, and the two scalars in which the sampling rate of \Cref{sec:main} is
expressed.

\begin{assumption}[GTR substitution model]\label{asm:gtr}
Sequences of length $\ell$ over a $d$-character alphabet evolve along $\mathcal{T}$ under a General
Time-Reversible rate matrix $Q$, as in \Cref{sec:bg-model}: conditions (i)--(iv) hold, and
consequently $\mathrm{tr}(Q)<0$.
\end{assumption}

\begin{assumption}[Molecular clock with contemporaneous terminal nodes]\label{asm:clock}
Branch lengths measure evolutionary time: the substitution process along an edge is determined by
$Q$ and by the time difference between its endpoints alone, with no per-lineage rate variation. All
$m$ terminal nodes are sampled at the same time, at height $\tau_{x}=0$.
\end{assumption}

Under \Cref{asm:clock} every terminal node sits at height $0$, so the path from a terminal node up
to any ancestor $h$ has length $\tau_h$. Both legs of the path through an MRCA are then equal, the
tree is \emph{ultrametric}, and the path length of \eqref{eq:mrca_height} is
\begin{equation}\label{eq:ultrametric}
   d_{ij} = 2\,\tau_{ij}.
\end{equation}
The closed form \eqref{eq:similarity_gtr} accordingly reads
\begin{equation}\label{eq:similarity_clock}
    S_{ij} = \exp\!\bigl(2\,\mathrm{tr}(Q)\,\tau_{ij}\bigr),
\end{equation}
so $S$ records the coalescence heights of $\mathcal{T}$, one entry per pair.

That is what places $S$ in the model on which spectral bipartitioning is analysed: the
\emph{constant block model} of Balakrishnan et al.\ \cite{Balakrishnan2011clustering}, in the
localized form used for latent trees by \cite{aizenbud2023spectral}.

\begin{definition}[Constant block model at a bipartition]\label{def:cfbm}
A symmetric matrix $S$ with unit diagonal satisfies the \emph{constant block model} at a
bipartition $\mathcal{X} = A \cup B$ if the cross-clan similarities are constant at a baseline
$\bze \in (0,1)$ while the intra-clan similarities clear that baseline from above by a floor,
\begin{equation}\label{eq:cfbm_floor}
   S_{ij}=\bze\quad(i\in A,\ j\in B),
   \qquad
   \bze<\Sin^{\min}\le S_{ij}\le\Smax<1\quad(i\ne j\ \text{in the same clan}),
\end{equation}
where $\Sin^{\min}$ is the weakest intra-clan similarity and $\Smax := \max_{i \ne j} S_{ij}$. The
diagonal blocks are otherwise arbitrary.
\end{definition}

\noindent An arbitrary clan pair of $\mathcal{T}$ does not satisfy \eqref{eq:cfbm_floor}: the pairs
it separates have MRCAs at many different heights, so \eqref{eq:similarity_clock} spreads the
cross-clan entries over a range. The root split is the exception. Every pair it separates has the
root itself as its most recent common ancestor, so \eqref{eq:similarity_clock} evaluates all of
those pairs at the single height $\tau_r$ and the cross-clan block collapses to one value.

\begin{proposition}[The population similarity is a constant block model at the root split]\label{prop:cfbm}
Let $S$ be the population similarity matrix of \Cref{def:simgraph} under \Cref{asm:gtr,asm:clock},
and suppose every branch of $\mathcal{T}$ has positive length. Then $S$ satisfies \Cref{def:cfbm} at
the root split $\mathcal{X} = A \cup B$, with
\begin{equation}\label{eq:cfbm_from_model}
   \bze = \exp\bigl(2\,\mathrm{tr}(Q)\,\tau_r\bigr),
   \qquad
   \Sin^{\min} = \exp\bigl(2\,\mathrm{tr}(Q)\,\tau^{\max}\bigr) > \bze,
\end{equation}
where $\tau_r$ is the root height and $\tau^{\max} < \tau_r$ is the greatest within-clan MRCA
height. In particular $S_{ii}=1$ and $0 < \bze \le S_{ij} \le \Smax < 1$ for every $i \ne j$, and the
block degrees of $L = D - S$ are
$D_A=\operatorname{diag}(S_A\1_{|A|})+|B|\bze I_{|A|}$ and
$D_B=\operatorname{diag}(S_B\1_{|B|})+|A|\bze I_{|B|}$.
\end{proposition}

\noindent \Cref{app:comp1} proves it. With the terminal nodes ordered along the split as in
\Cref{sec:bg-model}, the matrix and its Laplacian take the block form
\begin{equation}\label{eq:block_forms}
   S=\begin{pmatrix} S_A & \bze J_{|A|\times|B|}\\ \bze J_{|B|\times|A|} & S_B\end{pmatrix},
   \qquad
   L=\begin{pmatrix} D_A-S_A & -\bze J_{|A|\times|B|}\\ -\bze J_{|B|\times|A|} & D_B-S_B\end{pmatrix},
\end{equation}
with $J$ the all-ones block. The diagonal blocks are whatever the two subtrees make them; only the
off-diagonal block is pinned, and pinning it is what the model is for --- it is the single structural
fact the spectral analysis of \Cref{sec:bridging} runs on.

Two scalars carry the split through the analysis: how unevenly it divides the terminal nodes, and
how sharply intra-clan similarities exceed cross-clan ones.

\begin{definition}[Imbalance ratio and structural margin]\label{def:imbalance}
For clan sizes $|A| \le |B|$ with $|A|+|B|=m$, the \emph{imbalance ratio} is
\begin{equation}\label{eq:imbalance_ratio}
   \eta:=\frac{|B|}{|A|}\ \ge\ 1 .
\end{equation}
Writing $\Sout^{\max}$ and $\Sout^{\min}$ for the largest and smallest similarity across the split,
and $\Sin^{\min}$, $\Smax$ as in \eqref{eq:cfbm_floor}, the \emph{structural margin} is
\begin{equation}\label{eq:margin_definition}
   \rho := (\Sin^{\min}-\Sout^{\max})-(\Sout^{\max}-\Sout^{\min}).
\end{equation}
\end{definition}

\noindent A balanced split has $\eta=1$, while caterpillar-like lineages reach
$\eta=\mathcal{O}(m)$; $\eta$ plays the role of the balanced-cluster constant of Balakrishnan et
al.\ \cite{Balakrishnan2011clustering}. In $\rho$, the first bracket separates the weakest
within-clan pair from the strongest cross-clan pair and the second penalizes variation inside the
cross-clan block, following the range-restriction framework of the same work. Under
\Cref{asm:clock} that variation vanishes, $\Sout^{\max}=\Sout^{\min}=\bze$ by \Cref{prop:cfbm}, and
the penalty collapses:
\begin{equation}\label{eq:margin_under_clock}
   \rho=\Sin^{\min}-\bze > 0 .
\end{equation}

\Cref{fig:fiedler_partitions} shows the Fiedler split on three simulated trees spanning the range of
$\eta$: the smaller the clan the rule cuts off, the larger the values its terminal nodes carry.

\begin{figure}[t]
\centering
\begin{subfigure}[t]{0.30\textwidth}
  \captionsetup{singlelinecheck=false, justification=raggedright, skip=1pt}
  \subcaption{$\eta=1.5$}\label{fig:fied-eta01}
  \includegraphics[width=\textwidth]{Fiedler_Bipartitions/fiedler_tree_partition_eta01.png}
\end{subfigure}
\hfill{\color{black!25}\rule[-5.0cm]{0.4pt}{5.0cm}}\hfill
\begin{subfigure}[t]{0.30\textwidth}
  \captionsetup{singlelinecheck=false, justification=raggedright, skip=1pt}
  \subcaption{$\eta=4.4$}\label{fig:fied-eta05}
  \includegraphics[width=\textwidth]{Fiedler_Bipartitions/fiedler_tree_partition_eta05.png}
\end{subfigure}
\hfill{\color{black!25}\rule[-5.0cm]{0.4pt}{5.0cm}}\hfill
\begin{subfigure}[t]{0.30\textwidth}
  \captionsetup{singlelinecheck=false, justification=raggedright, skip=1pt}
  \subcaption{$\eta=15.1$}\label{fig:fied-eta15}
  \includegraphics[width=\textwidth]{Fiedler_Bipartitions/fiedler_tree_partition_eta15.png}
\end{subfigure}
\caption{$\Pisp{S}$ on three trees, one per panel. Each is drawn
from a pool of simulated Kingman trees ($m=500$ terminal nodes, JC69 sequences of length
$\ell=10^4$) binned by the imbalance its Fiedler split realises. In each panel the
tree is on top, its branches coloured by the sign of the Fiedler vector --- grey on
the one branch that spans both signs --- and below it the Fiedler value of each terminal node
is plotted at that node's own horizontal position in the tree. The sign pattern cuts a single edge
every time, as \Cref{thm:stdr-split} requires, and the smaller the clan it cuts off, the
larger the values carried by its terminal nodes.}
\label{fig:fiedler_partitions}
\end{figure}

\Cref{asm:gtr,asm:clock} are hypotheses on how the data are generated, and they already fix every
quantity in \eqref{eq:margin_definition}. What they do not fix is whether those quantities stand in
the relation the analysis needs. That relation is the following inequality. It is not an assumption
about the data but a property checkable on any given tree, and it is what delimits the trees the
guarantee covers; it is the only condition beyond \Cref{asm:gtr,asm:clock} that the main result
carries.

\begin{condition}[Separation margin]\label{asm:margin}
The structural margin strictly dominates the imbalance-weighted cross-clan boundary,
\begin{equation}\label{eq:margin_assumption}
   \rho>\eta\,\Sout^{\max},\qquad\text{equivalently}\qquad
   \Sin^{\min}>(1+\eta)\,\Sout^{\max},
\end{equation}
which under \Cref{asm:clock} reads $\Sin^{\min}>(1+\eta)\bze$.
\end{condition}
